\scanpage{160}
% archive foliation: 159
On ne suppose plus donné de $\unsure{\uline L|E}$

$A'')$ Équivalence des deux données suivantes

(i) Une $S$-classe $\xi$ de faisceaux inversibles sur $S$,\inserted{qui est \uline{\unsure{amplissime}}} et un isom
\[
\varphi:\mathrm P^N_{\unsure{S|\xi}}\xrightarrow{\sim}P_\xi
\]
(ii) Une Donnée d'un morphisme
\[
j:X\longrightarrow\mathrm P_S^N
\]
tel que \deleted{$\unsure{j^*(\uline L)}$}
\[
q_*(\uline O(1))\longrightarrow f_*(j^*(\uline O(1)))
\]
soit un isomorphisme
\margin{\uline{morphisme spécial}\quad ou \unsure{projectif}}
\margin{Attention, $2$\quad\inserted{spécial,}
\diagram{\begin{tikzcd}[row sep=large]
\uline L\ \text{\unsure{magnifique|amplissime} sur }X/S+\text{trivialisation projective}
\arrow[d]\arrow[ddd,bend right=65,"\unsure{isom induit par un homoth. identique}"']\\
S\text{-classe }\xi\text{ sur }S,+\text{trivialisation projective}\arrow[d,equal]\\
\text{Morphisme spécial }X-\mathrm P_S^N\arrow[d]\\
\uline L'\ \text{\unsure{magnifique|amplissime} spécial sur }X,\text{ et trivialisation projective}
\end{tikzcd}}
$\#$
}
\deleted{\unsure{Les données}\illegible{2}}

$A''')$ Équivalence entre les données

(i) Une $S$-classe de faisceaux sur $X$ \deleted{\unsure{très ample}\illegible{1}\ ``spécial''}\inserted{\unsure{amplissime} et spécial}

(ii) Un fibré de Brauer-Severi sur $S$, associé à un fibré \uline{principal} $\mathcal P$ de groupe $GP\deleted{\unsure{L}}(N)$
$+$ morphisme \uline{spécial}
\[
j:X_{\unsure{\mathcal P}}-P^N_{\unsure{\mathcal P}\deleted{\unsure{\mathrm{Pic}}}}
\]
\deleted{\unsure{le morphisme}}\inserted{ce morphisme étant \uline{covariant} pour $GP(N)$}

\margin{spécial $=$ \unsure{pour ch. fibre en dim} $0$, et les \unsure{points définissent un morphisme, vers un projectif entier, comme} \illegible{2} de \deleted{\unsure{dim.}}\inserted{\unsure{dim.}}\unsure{fibre} $N$}

Si on introduit $\unsure{\mathfrak J}\ \uuline{\mathrm{Hom}}_S(X,P^N)$, où $GP(N)$ opère,
le sous-schém. $\unsure{\mathcal J}:\uuline{\mathrm{Hom}}_S^{\mathrm{spéc}}(X,P^N)$
[cas $X$ \deleted{\unsure{propre}} projectif sur $S$ loc noeth], alors on voit que
$\deleted{GP}(N)$ opère \uline{sans pts fixes} sur $\uuline{\mathrm{Hom}}_S^{\mathrm{spéc}}(X,P^N)$,
et les structures $A'''$ sont

\margin{\unsure{pour les points géométriques}\illegible{3}\ \unsure{introduire} \uline{\unsure{projectif}}\ \unsure{de passage}\illegible{2}.
\unsure{le morphisme}\illegible{2}\ $\unsure{j^*(\uline O(1))}$ \unsure{est}\illegible{2}
\deleted{\unsure{fibre}} \unsure{sur les fibres}\illegible{2}\ $\unsure{\mu}$
\deleted{$\unsure{\mathrm{Hom}_{\mathcal P}(X,\mathrm P^N)}$}\ $\unsure{\mathrm{Hom}}$}

